\section{Decomposition theorem in the quantum setting}
\label{section:quantum-setting}


\subsection{Finite W-algebra for the minimal orbit}

Finite W-algebras are defined by a quantum Hamiltonian reduction process \cite{kostant1978whittaker,premet2002special,gan2002quantization,brundan2007good}. See \cite[Section 7.1]{gan2006almost} for generalities on quantum Hamiltonian reduction.

Let $\U(\g)$ be the universal enveloping algebra of $\g$. It is equipped with an adjoint action of $G$ by algebra automorphisms and an induced adjoint action of $\g$. The map
\[ \widetilde{\mu} \colon \g_{\geqslant 1} \longrightarrow \U(\g), \quad a \longmapsto a + \chi(a) 1 \]
is a \emph{quantum comoment map} for the $G_{\geqslant 1}$-action, meaning that it is a Lie algebra homomorphism such that the following equality holds:
\[ [\widetilde{\mu}(a), F] = a \cdot F, \quad \text{for} \quad a \in \g_{\geqslant 1} \quad \text{and} \quad F \in \U(\g), \]
where $a \cdot F$ denotes the adjoint action.

The \emph{minimal finite W-algebra} is the quantum Hamiltonian reduction
\[ U(\g, f) \coloneqq (\U(\g) / I_{\widetilde{\mu}})^{G_{\geqslant 1}}, \]
where $I_{\widetilde \mu}$ is the left ideal of $\U(\g)$ that is generated by the elements $x_{1,2} + 1$ and $x_{1,i}$ for $3 \leqslant i \leqslant d$, that is to say by the image of $\widetilde{\mu}$. By a standard fact on quantum Hamiltonian reduction, $\U(\g, f)$ inherits an associative algebra structure from $\U(\g)$.


\subsection{The Kazhdan filtration on the W-algebra} We recall the definition of the Kazhdan filtration on $\U(\g)$ and $\U(\g, f)$ from \cite[Section 4]{gan2002quantization}. The semisimple element $H$ acts on $\g$ by derivations and its action is diagonalisable, defining the good grading. Therefore, we get a well-defined diagonalisable action on $\U(\g)$. The \emph{Kazhdan filtration} is defined by
\[ \U(\g)_p \coloneqq \sum_{n - m \leqslant p} \operatorname{Span} \{ x_1 \cdots x_n \in \U(\g) \quad | \quad [H, x_1 \cdots x_n] = m \, x_1 \cdots x_n \}, \]
where $p$ belongs to $\Z$.

Because $\g = \g_{-1} \oplus \g_0 \oplus \g_1$ and $e$ is in $\g_1$, $e$ is in the level $0$ for this filtration and $\U(\g)_p$ is non-zero only if $p \geqslant 0$. 

\begin{remark} \label{remark:nonnegative-filtration}
	The fact that the filtration is supported on $\Z_{\geqslant 0}$ comes from the fact that the good grading on $\g$ is supported on $\{-1, 0, 1\}$. If $f$ is not a minimal nilpotent element, it is almost always the case that this Kazhdan grading contains components with negative degrees. For example, take the regular nilpotent element in $\sl_3$.
\end{remark}

The Kazhdan filtration induces filtrations on the quotient $\U(\g) / I_{\widetilde{\mu}}$ and on the finite W-algebra $\U(\g, f) = (\U(\g) / I_{\widetilde{\mu}})^{G_{\geqslant 1}}$. They are both called \emph{Kazhdan filtrations} too. The following properties are easy to prove, see \cite{gan2002quantization} or \cite[Section 4.2.1]{juillard2025phd}.

\begin{lemma}
	\begin{enumerate}
		\item The Kazhdan filtration is exhaustive.
		
		\item It satisfies the multiplicative conditions \eqref{equation:multiplicative-properties}, so $\gr \U(\g)$ is commutative.
		
		\item There is an isomorphism $\C[\g^\vee] = \Sym \g \cong \gr \U(\g)$. 
		
		\item The commutative subalgebra $\U(\g)_0$ is a free commutative algebra generated by the subspace $\g_1$:
		\[ \U(\g)_0 \cong \Sym \g_1. \]
	\end{enumerate}
\end{lemma} 

Because $e$ belongs to $\U(\g)_0$, we can apply the localisation procedure described in \Cref{subsectin:localisation-filtered}. Then, $\U(\g)[e^{-1}]$ is well-defined, it inherits a Kazhdan filtration and there is a Poisson algebra isomorphism
\[ \C[\g^\vee[e^{-1}]] = (\Sym \g)[e^{-1}] \cong \gr \big(\U(\g)[e^{-1}]\big). \]

The following theorem proves that the Slodowy slice is quantised by the finite W-algebra.

\begin{theorem}[{\cite[Theorem 4.1]{gan2002quantization}}] \label{theorem:slodowy-quantisation}
	The isomorphism $\Sym \g \cong \gr \U(\g)$ induces algebra isomorphisms
	\[ \Sym \g / I_\mu \cong \gr(\U(\g) / I_{\widetilde{\mu}}) \quad \text{and} \quad \C[S_f] \cong (\Sym \g / I_\mu)^{G_{\geqslant 1}} \cong \gr \U(\g, f), \]
	and the second isomorphism is Poisson.
\end{theorem}



\subsection{The W-algebra is a subalgebra of the localisation}

We state a quantisation of \Cref{theorem:slodowy-git}, that is the reconstruction of the finite W-algebra as a subalgebra of $\U(\g)[e^{-1}]$. Because the punctured line $\C^\times e$ is stable under the adjoint action of $K \ltimes G_{\geqslant 1}$, the action of this group extends from $\U(\g)$ to the localisation $\U(\g)[e^{-1}]$.

If we consider the opposite grading $\g = \bigoplus_{i \in \Z}\g_{-i}$, then the Lie subalgebra $\kfk \ltimes \g_{\geqslant 1}$ is $\Z_{\leqslant 0}$-graded and $\U(\g)[e^{-1}]$ is a filtered module because the Kazhdan filtration satisfies the property:
\[ [x, \big(\U(\g)[e^{-1}]\big)_p] \subseteq \big(\U(\g)[e^{-1}]\big)_{1 - i + p - 1} = \big(\U(\g)[e^{-1}]\big)_{p - i}, \]
for $x$ in $\g_i \subseteq \U(\g)_{1 - i}$ and $p$ an integer.

\begin{theorem} \label{theorem:finite-w-algebra-coset}
	There is an algebra isomorphism
	\[ \big(\U(\g)[e^{-1}]\big)^{K \ltimes G_{\geqslant 1}} \cong \U(\g, f).  \]
\end{theorem}

To prove this theorem, we start by computing the associated graded algebra of $\big(\U(\g)[e^{-1}]\big)^{K \ltimes G_{\geqslant 1}}$ for the induced Kazhdan filtration.

\begin{proposition} \label{proposition:gr-invariant-localisation}
	The natural algebra homomorphism
	\[ \gr \big(\big(\U(\g)[e^{-1}]\big)^{K \ltimes G_{\geqslant 1}}\big) \longrightarrow \big(\gr \big(\U(\g)[e^{-1}]\big)\big)^{K \ltimes G_{\geqslant 1}} \]
	is an isomorphism. Therefore, we have a Poisson algebra isomorphism
	\[ \gr \big(\big(\U(\g)[e^{-1}]\big)^{K \ltimes G_{\geqslant 1}}\big) \cong \C[\g^\vee[e^{-1}] \git (K \ltimes G_{\geqslant 1})]. \]
\end{proposition}

\begin{proof}
	The group $K = \Gm \times \Ga{}^{d-2}$ decomposes as the product of a torus and a unipotent group.
	
	To prove the proposition, we can first deal with the action of the unipotent normal subgroup $\Ga{}^{d-2} \ltimes G_{\geqslant 1}$ in $K \ltimes G_{\geqslant 1}$. To show the isomorphism
	\begin{equation} \label{equation:gr-invariant-unipotent} \gr \big(\big(\U(\g)[e^{-1}]\big)^{\Ga{}^{d-2} \ltimes G_{\geqslant 1}}\big) \cong \big(\gr \big(\U(\g)[e^{-1}]\big)\big)^{\Ga{}^{d-2}\ltimes G_{\geqslant 1}}, \end{equation}
	we want to use \Cref{proposition:gr-invariant}. 
	
	According to \Cref{proposition:dilation-gan-ginzburg}, we have the $\Ga{}^{d-2} \ltimes G_{\geqslant 1}$-equivariant isomorphism
	\[ \C[\Ga{}^{d-2} \ltimes G_{\geqslant 1}] \otimes \C[\Gm \times S_f] \cong (\Sym \g)[e^{-1}] \cong \gr (\U(\g)[e^{-1}]).  \]
	With the discussion on the Lie algebra grading at the beginning of the section, we have all the assumptions to apply \Cref{proposition:gr-invariant}, and we get the isomorphism~\eqref{equation:gr-invariant-unipotent}.
	
	The isomorphism \eqref{equation:gr-invariant-unipotent} is $\Gm$-equivariant, so we get the isomorphism
	\[ \big(\gr \big(\big(\U(\g)[e^{-1}]\big)^{\Ga{}^{d-2} \ltimes G_{\geqslant 1}}\big)\big)^{\Gm} \cong \big(\gr \big(\U(\g)[e^{-1}]\big)\big)^{K \ltimes G_{\geqslant 1}}. \]
	between the $\Gm$-invariant subalgebras. By applying \Cref{proposition:weight-filtration} on the left-hand side, we get the desired isomorphism:
	\[ \gr \big(\big(\U(\g)[e^{-1}]\big)^{K \ltimes G_{\geqslant 1}}\big) \cong \big(\gr \big(\U(\g)[e^{-1}]\big)\big)^{K \ltimes G_{\geqslant 1}}. \]
	
	The second isomorphism of the proposition follows because we have the Poisson isomorphism
	\[ \gr \big(\U(\g)[e^{-1}]\big) \cong (\Sym \g)[e^{-1}] = \C[\g^\vee[e^{-1}]] \]
	which is $K \ltimes G_{\geqslant 1}$-equivariant. By taking the invariant subalgebras, we get:
	\[ \gr \big(\big(\U(\g)[e^{-1}]\big)^{K \ltimes G_{\geqslant 1}}\big) \cong \big(\gr \big(\U(\g)[e^{-1}]\big)\big)^{K \ltimes G_{\geqslant 1}} \cong \C[\g^\vee[e^{-1}] \git (K \ltimes G_{\geqslant 1})]. \]
\end{proof}

Now, we can prove the main theorem.

\begin{proof}[Proof of \Cref{theorem:finite-w-algebra-coset}]
	First, the localisation homomorphism $\U(\g) \hookrightarrow \U(\g)[e^{-1}]$ induces an algebra isomorphism
	\[ (\U(\g) / I_{\widetilde \mu} )^{G_{\geqslant 1}} \cong (\U(\g)[e^{-1}] / I_{\widetilde \mu} [e^{-1}])^{G_{\geqslant 1}}. \]
	The proof is similar to that of \Cref{proposition:hamiltonian-reduction-open}: we use the exactness of localisation and the fact that $e$ is already invertible in the quotient $\U(\g) / I_{\widetilde \mu}$, so we have an isomorphism
	\[ \U(\g) / I_{\widetilde \mu}  \cong \U(\g)[e^{-1}] / I_{\widetilde \mu} [e^{-1}] \]
	according to \cite[Proposition 3.6]{stenstrom1975rings}.
	
	Second, there is a natural algebra homomorphism
	\[ \lambda\colon (\U(\g)[e^{-1}])^{K \ltimes G_{\geqslant 1}} \longrightarrow  (\U(\g)[e^{-1}] / I_{\widetilde \mu} [e^{-1}])^{G_{\geqslant 1}} \cong \U(\g, f). \]
	Both sides are equipped with Kazhdan filtrations that are bounded from below and exhaustive. So if $\gr \lambda$ is an isomorphism, then \Cref{lemma:quantisation-iso} implies that $\lambda$ is an isomorphism.
	
	By \Cref{proposition:gr-invariant-localisation} and \Cref{theorem:slodowy-quantisation}, we have the algebra isomorphisms
	\[ \begin{multlined}
		\gr \big(\big(\U(\g)[e^{-1}]\big)^{K \ltimes G_{\geqslant 1}}\big) \cong \C[\g^\vee[e^{-1}] \git (K \ltimes G_{\geqslant 1})] \\ \text{and} \quad \gr \U(\g, f) \cong \C[\mu^{-1}(0) \git G_{\geqslant 1}].
	\end{multlined} \]
	Under these isomorphisms, $\gr \lambda$ is conjugate to the isomorphism
	\[ \C[\g^\vee[e^{-1}] \git (K \ltimes G_{\geqslant 1})] \cong  \C[\mu^{-1}(0) \git G_{\geqslant 1}] \]
	given by \Cref{theorem:slodowy-git}. Therefore, $\gr \lambda$, and then $\lambda$, are isomorphisms.	
\end{proof}


\subsection{Decomposition theorem for finite W-algebras}

\begin{lemma}
	The two subalgebras
	\[ \U(\kfk \ltimes \g_{\geqslant 1})[e^{-1}] \quad \text{and} \quad \big(\U(\g)[e^{-1}]\big)^{K \ltimes G_{\geqslant 1}} \cong \U(\g, f)  \]
	of $\U(\g)[e^{-1}]$ commute with each other.
\end{lemma}

\begin{proof}
	The lemma follows immediately from the following equalities:
	\[ \begin{aligned}
		\big(\U(\g)[e^{-1}]\big)^{K \ltimes G_{\geqslant 1}} & = \big(\U(\g)[e^{-1}]\big)^{\kfk \ltimes \g_{\geqslant 1}} \\
		& = \big\{ F \in \U(\g)[e^{-1}] \quad | \quad \text{for all} ~  x \in \kfk \ltimes \g_{\geqslant 1},  ~   [x, F] = 0 \big\} \\
		& = \big\{ F \in \U(\g)[e^{-1}] \quad | \quad  [\U(\kfk \ltimes \g_{\geqslant 1})[e^{-1}], F]  = 0\big\}. 
	\end{aligned} \]
	\begin{itemize}
		\item The first equality is justified by the connectedness of the group $K \ltimes G_{\geqslant 1}$. 
		\item The second equality is because the bracket $[x, \bullet]$ coincides with the infinitesimal action of $x$ in $\kfk \ltimes \g_{\geqslant 1}$.
		
		\item The third equality follows because commuting with $\U(\kfk \ltimes \g_{\geqslant 1})[e^{-1}]$ is equivalent to commuting with $\kfk \ltimes \g_{\geqslant 1}$, since this space generates the non-localised algebra and the bracket with $e^{-1}$ is given by the formula: $[e^{-1}, \bullet] = - e^{-1} [e, \bullet] e^{-1}$.
	\end{itemize}
\end{proof}

Consequently, there is a natural algebra homomorphism
\[ \widetilde{\delta} \colon \U(\kfk \ltimes \g_{\geqslant 1})[e^{-1}] \otimes \big(\U(\g)[e^{-1}]\big)^{K \ltimes G_{\geqslant 1}} \longrightarrow \U(\g)[e^{-1}] \]
given by multiplication. This provides the decomposition isomorphism for the finite W-algebra.

\begin{theorem} \label{theorem:decomposition-quantum}
	The algebra homomorphism
	\[ \widetilde{\delta} \colon \U(\kfk \ltimes \g_{\geqslant 1})[e^{-1}] \otimes \big(\U(\g)[e^{-1}]\big)^{K \ltimes G_{\geqslant 1}} \longrightarrow \U(\g)[e^{-1}] \]
	is an isomorphism. Therefore, we have an algebra isomorphism
	\[ \U(\kfk \ltimes \g_{\geqslant 1})[e^{-1}] \otimes \U(\g, f) \cong \U(\g)[e^{-1}]. \]
\end{theorem}

\begin{proof}
	Both sides are equipped with induced filtrations coming from the Kazhdan filtration on $\U(\g)$ and localising at $e$ in $\U(\kfk \ltimes \g_{\geqslant 1})_0 \subseteq \U(\g)_0$. Because the filtrations on both sides are bounded from below and exhaustive, we just need to prove that
	\[ \gr \widetilde\delta \colon \gr \big(\U(\kfk \ltimes \g_{\geqslant 1})[e^{-1}]\big) \otimes \gr \big(\big(\U(\g)[e^{-1}]\big)^{K \ltimes G_{\geqslant 1}}\big) \longrightarrow \gr\big( \U(\g)[e^{-1}]\big) \]
	is an isomorphism to apply \Cref{lemma:quantisation-iso}.
	
	By applying \Cref{proposition:gr-localisation}, we obtain the algebra isomorphisms:
	\[ \gr\big(U(\kfk \ltimes \g_{\geqslant 1}\big)[e^{-1}]) \cong \Sym(\kfk \ltimes \g_{\geqslant 1})[e^{-1}]  \]
	and 
	\[ \gr\big(\U(\g)[e^{-1}]\big) \cong \Sym(\g)[e^{-1}]. \]
	
	Moreover, \Cref{proposition:gr-invariant-localisation} gives the algebra isomorphism
	\[ \gr \big(\big(\U(\g)[e^{-1}]\big)^{K \ltimes G_{\geqslant 1}}\big) \cong \big((\Sym\g)[e^{-1}]\big)^{K \ltimes G_{\geqslant 1}}.  \]
	
	Under these isomorphisms, $\gr \widetilde{\delta}$ is conjugate to the Poisson algebra isomorphism 
	\[ \delta^* \colon \Sym(\kfk \ltimes \g_{\geqslant 1})[e^{-1}] \otimes \big((\Sym \g)[e^{-1}]\big)^{K \ltimes G_{\geqslant 1}}  \overset{\sim}{\longrightarrow} (\Sym \g)[e^{-1}]  \]
	provided by \Cref{theorem:decomposition-slodowy}. This concludes the proof.
\end{proof}
